
\section{Definitions}
\label{app:defs}

\paragraph{Labels alone.} For $s$ positives among $n$ 0/1 labels, Clopper-Pearson's upper limit is
$\mathrm{Beta}^{-1}(1-\delta;\,s+1,\,n-s)$, and $1-\delta^{1/n}$ at $s=0$. With no positive, the smallest sample that
certifies a threshold $\tau$ is $n\ge\ln\delta/\ln(1-\tau)$: 299 labels for 1\% at $\delta=0.05$, 149 for 2\%, 59 for
5\%. For labels in $[0,1]$ the exact choices are Hoeffding's \citep{hoeffding1963}, Bentkus's \citep{bentkus2004},
Anderson's \citep{learnedmiller2019} and the betting bound \citep{waudbysmith2024betting}. Wilson's bound
\citep{wilson1927} is the smallest $m>\hat p$ with $m-\hat p\ge z_{1-\delta}\sqrt{m(1-m)/n}$.

\paragraph{Stratified limits.} With strata $h$ of population share $W_h$, $n_h$ labels, $s_h$ positives and rates
$\hat p_h$, $\hat\mu=\sum_hW_h\hat p_h$. The \SWL{} is the smallest $m>\hat\mu$ with
$m-\hat\mu\ge z_{1-\delta}\sqrt{\sum_hW_h^2p_h(m)(1-p_h(m))/n_h}$, where $p_h(m)$ shifts every stratum's rate by
$m-\hat\mu$, clipped to $[0,1]$; it is a heuristic, not the inversion of a stratified score test. The \SWT{} is
$\hat\mu+t\sqrt{\sum_hW_h^2\tilde p_h(1-\tilde p_h)/(n_h-1)}$ with $\tilde p_h=(s_h+1/2)/(n_h+1)$ and a Student-t
quantile at Satterthwaite's degrees of freedom \citep{satterthwaite1946}. For a claim about the population a pool of $N$
prompts was drawn from, both add $\sum_hW_h(\hat p_h-\hat\mu)^2/N$ to the variance, the first-phase term of two-phase
sampling \citep{neyman1938,cochran1977}.

\paragraph{A judge.} With $n$ gold labels $e$, the judge's feature $f$ on them and on $N_u$ unlabelled responses, PPI++
\citep{angelopoulos2023ppipp} estimates $\hat p=\overline{e}+\gamma(\overline{f}_{\text{unl}}-\overline{f}_{\text{lab}})$ with
$\gamma=\mathrm{cov}(e,f)/((1+n/N_u)\mathrm{var}(f))$ and variance $V=\mathrm{var}(e-\gamma f)/n+\gamma^2\mathrm{var}_{\text{unl}}(f)/N_u$;
the published limit is $\hat p+z_{1-\delta}\sqrt V$. StratPPI \citep{fisch2024stratppi} applies the correction within
strata. PPBoot \citep{zrnic2024ppboot} uses a percentile of bootstrap estimates. The bootstrap-t limit
\citep{hall1988} resamples the labelled responses $B$ times (within strata where stratified), redraws each unlabelled
mean, and sets $U=\hat p-t^*_\delta\sqrt V$ with $t^*_\delta$ the $\delta$-quantile of $(\hat p^*_b-\hat p)/\sqrt{V^*_b}$.

\paragraph{Clustered data.} With clusters $c$ (user tasks) of sizes $n_c$ and successes $s_c$, $\hat p=\sum s_c/\sum
n_c$ and $se=\sqrt{C/(C-1)\sum_c(s_c-\hat pn_c)^2}/\sum_cn_c$; the limit is a bootstrap-t with whole clusters resampled.
For crossed clusters, the multiway variance is $V_u+V_{\text{inj}}-V_p$ \citep{cameron2011multiway}, and the pigeonhole
bootstrap \citep{owen2007pigeonhole} resamples both kinds of task independently.

\paragraph{Effective sample size.} $\mathrm{ESS}=(w_{\text{base}}/w_A)^2$ at equal labels and $\delta$. Unless stated,
$w$ is measured from the truth ($\mathrm{ESS}_{\text{truth}}$); the first strata study measures it from the estimate
($\mathrm{ESS}_{\text{est}}$). A draw that returns no limit enters $w$ as a limit of 1, so abstaining lowers a bound's ESS and cannot raise it. A design that removes a share $G$ of the variance has a predicted ESS of $1/(1-G)$; for
reference-rate strata $G\approx\mathrm{ICC}_{\text{cand}}\,\omega^2\,\mathrm{rel}(k)\,c_H$, with the updated policy's
intraclass correlation, the correlation $\omega$ of a prompt's pre- and post-update rates, the reliability
$\mathrm{rel}(k)=k\,\mathrm{ICC}_{\text{ref}}/(1+(k-1)\mathrm{ICC}_{\text{ref}})$ of a $k$-sample reference rate, and
$c_H=0.96$ at $H=8$. A judge with squared correlation $\rho^2$ to the gold label gives about $1/(1-\rho^2)$ when
unlabelled responses are plentiful.

\paragraph{The carried bound.} With recall $r$, false-alarm rate $f$ and flag rate $q=rp+f(1-p)$, $p=(q-f)/(r-f)$
\citep{rogan1978}. The carried bound is $U=(q^{+}-f^{-})/(r^{-}-f^{-})$, where $q^{+}$ is a Clopper-Pearson upper limit
on the flag rate of the certified population and $r^{-}$, $f^{-}$ are Clopper-Pearson lower limits on recall and the
false-alarm rate from the calibration labels, each at $\delta/3$; $p$ is increasing in $q$ and, for $f\le q\le r$, decreasing in
$r$ and $f$. It returns no certificate ($U=1$) when $r^{-}\le f^{-}$, when $q^{+}<f^{-}$ (which the model rules out), or when the
calibration holds too few positives to bound recall.

\section{The Full Validity Table}
\label{app:validity}

Tables~\ref{tab:fullA} to~\ref{tab:fullD} give every row. Rows in bold are bounds we do not use in that setting and
that are over their level there. Rows 1, 3 to 5 and 15 are exact computations over binomial counts on a harm pool from a
companion training run (Qwen2.5-0.5B policies; an unpublished paper) whose per-response labels were lost; its printed rate pins the count to
121 to 124 positives of 3,600, and every count in that range is computed exactly. Rows 37 to 55 were registered before they were run, and so were rows 14, 30 and 31, by a commit whose order
rests on its local time (Appendix~\ref{app:record}).

\begin{table*}[p]
\centering
\scriptsize
\setlength{\tabcolsep}{3pt}
\begin{tabular}{@{}r p{2.8cm} p{2.0cm} p{3.9cm} l p{2.1cm} p{2.0cm} p{1.9cm}@{}}
\toprule
\textbf{\#} & \textbf{bound} & \textbf{kind} & \textbf{setting} & \textbf{$\delta$} & \textbf{miss, per cell} & \textbf{cells over / unresolved / at or under} & \textbf{largest 95\% upper limit of a cell's miss (cells over after correction)} \\
\midrule
1 & Clopper-Pearson & exact (binary labels) & binomial counts at the rate of a trained-policy harm pool (labels lost), 121--124 positives of 3,600, n 200--2,400 & 0.10 & 0.066--0.096 & 0 / 0 / 20 & 0.096 (0) \\
2 & Clopper-Pearson & exact & the judge plasmodes (App.~\ref{app:judge}), every cell & 0.05 & 0--0.052 & 0 / 2 / 40 & 0.059 (0) \\
3 & betting mixture & exact (bounded) & same pool as row 1 & 0.10 & 0.007--0.018 & 0 / 0 / 20 & 0.018 (0) \\
4 & Bentkus & exact (bounded) & same pool & 0.10 & 0.017--0.037 & 0 / 0 / 20 & 0.037 (0) \\
5 & Hoeffding, Anderson & exact (bounded) & same pool & 0.10 & 0.000 & 0 / 0 / 40 & 0.000 (0) \\
6 & stratified Wilson-type limit & approximate & 4 mid-rate labels (9--95\%), real Granite-3.3-2B responses, $n_s$ 100--200, 3 checkpoints & 0.05 / 0.10 & 0.017--0.054 / 0.058--0.097 & 0 / 3 / 21; 0 / 0 / 24 & 0.061 (0); 0.106 (0) \\
7 & the Wilson-type limit & approximate & label pushed by the Lagrangian, $n_s$ 200 & 0.05 / 0.10 & 0.011--0.023 / 0.040--0.064 & 0 / 0 / 2; 0 / 0 / 2 & 0.027 (0); 0.071 (0) \\
8 & the Wilson-type limit on a design-weighted sheet & approximate & sheets re-drawn by their real sampling rule & 0.05 & 0--0.060 & 1 / 1 / 16 & 0.067 (0) \\
9 & PPI++ with a bootstrap-t limit & approximate & the judge plasmodes (App.~\ref{app:judge}), every cell and feature & 0.05 & 0--0.054 & 0 / 13 / 155 & 0.061 (0) \\
10 & cluster bootstrap-t, by user task & approximate & AgentDojo, 28 pipelines, scheme (a) & 0.05 & 0.013--0.051 & 0 / 2 / 26 & 0.058 (0) \\
11 & StratPPI estimator with a bootstrap-t limit, proportional allocation & approximate & reference-rate strata, 5 mid-rate labels, $n_s$ 100--200; judge-logit strata, n 100--1,000 & 0.05 & 0.024--0.045; 0.001--0.056 & 0 / 0 / 28; 0 / 5 / 23 & 0.051 (0); 0.064 (0) \\
12 & the same estimator and limit, redrawn with replacement & approximate & reference-rate strata, 5 mid-rate labels, $n_s$ 100--200 & 0.05 / 0.10 & 0.030--0.050 / 0.069--0.100 & 0 / 0 / 28; 0 / 1 / 27 & 0.052 (0); 0.103 (0) \\
13 & Clopper-Pearson, random draws with replacement (the control) & exact & the same 5 labels and sizes & 0.05 / 0.10 & 0.029--0.050 / 0.050--0.097 & 0 / 0 / 28; 0 / 0 / 28 & 0.052 (0); 0.100 (0) \\
14 & pigeonhole bootstrap-t & approximate & AgentDojo, 28 pipelines, schemes (a); (b); (c) & 0.05 & at most 0.004; 0.037; 0.042 & 0 / 0 / 28; 0 / 0 / 28; 0 / 0 / 28 & 0.007 (0); 0.043 (0); 0.048 (0) \\
15 & \textbf{Student-t} & fails at low rates & same pool as row 1, n 200--2,400 & 0.10 & \textbf{0.104--0.195} & 20 / 0 / 0 & 0.195 (20) \\
16 & \textbf{the Wilson-type limit and the pooled Wilson bound at rare rates} & in the stored design, not over at $\delta$ 0.05; over in 2 cells at 0.10 & harm labels at 1--2\%, $n_s$ 100--200, 3 checkpoints & 0.05 / 0.10 & 0--0.048 / \textbf{0--0.126} & 0 / 0 / 24; 2 / 1 / 21 & 0.054 (0); 0.136 (2) \\
17 & \textbf{the Wilson-type limit, redrawn with replacement} & over on the two labels at rates of 65\% and above & reference-rate strata, 5 mid-rate labels, $n_s$ 100--200 & 0.05 / 0.10 & \textbf{0.017--0.062 / 0.056--0.113} & 7 / 1 / 20; 6 / 2 / 20 & 0.065 (7); 0.116 (4) \\
18 & \textbf{pooled Wilson bound, random draws with replacement} & over in 6 and 10 of the 28 cells & the same 5 labels and sizes & 0.05 / 0.10 & \textbf{0.031--0.069 / 0.069--0.128} & 6 / 0 / 22; 10 / 3 / 15 & 0.071 (6); 0.131 (10) \\
\bottomrule
\end{tabular}

\caption{Miss rate against $\delta$, rows 1 to 18. Cells are counted as over / unresolved / at or under $\delta$; the next
column is the largest upper end of a 95\% interval for a cell's miss rate, with the cells still over after a
Bonferroni correction in brackets (for the exact binomial rows, the largest exact miss and the cells above $\delta$). Draws per cell: 4,000 to 40,000. AgentDojo schemes: (a) user tasks resampled, injection tasks fixed; (b) the reverse;
(c) both. ``Redrawn with replacement'' draws each stratum's prompts with replacement, the limit of a pool far larger
than the safety set.}
\label{tab:fullA}
\end{table*}

\begin{table*}[p]
\centering
\scriptsize
\setlength{\tabcolsep}{3pt}
\begin{tabular}{@{}r p{2.8cm} p{2.0cm} p{3.9cm} l p{2.1cm} p{2.0cm} p{1.9cm}@{}}
\toprule
\textbf{\#} & \textbf{bound} & \textbf{kind} & \textbf{setting} & \textbf{$\delta$} & \textbf{miss, per cell} & \textbf{cells over / unresolved / at or under} & \textbf{largest 95\% upper limit of a cell's miss (cells over after correction)} \\
\midrule
19 & \textbf{PPI++ with a normal limit} & fails in most cells & the judge plasmodes (App.~\ref{app:judge}) & 0.05 & \textbf{0.041--0.241} & 126 / 29 / 13 & 0.255 (87) \\
20 & \textbf{StratPPI's normal limit, proportional allocation} & fails at these sizes & reference-rate strata; judge-logit strata & 0.05 & \textbf{0.024--0.086; 0.054--0.239} & 9 / 5 / 14; 26 / 2 / 0 & 0.094 (7); 0.252 (25) \\
21 & \textbf{StratPPI's normal limit, the paper's allocations (oracle; heuristic)} & fails & reference-rate strata; judge-logit strata & 0.05 & \textbf{up to 0.110; 0.217; up to 0.219; 0.90} & 9 / 3 / 14; 11 / 0 / 15; 28 / 0 / 0; 24 / 1 / 3 & 0.119 (8); 0.229 (11); 0.232 (25); 0.912 (24) \\
22 & \textbf{StratPPI estimator with a bootstrap-t limit, oracle or heuristic allocation} & fails & the same cells & 0.05 & \textbf{up to 0.078; 0.174; up to 0.083; 0.90} & 6 / 3 / 17; 11 / 1 / 14; 14 / 11 / 3; 14 / 0 / 14 & 0.086 (5); 0.185 (11); 0.092 (9); 0.910 (14) \\
23 & \textbf{PPBoot, percentile limit (basic; power-tuned)} & fails at these sizes & judge-logit cells, unstratified & 0.05 & \textbf{up to 0.133; 0.148} & 7 / 6 / 1; 12 / 2 / 0 & 0.144 (5); 0.160 (11) \\
24 & \textbf{Clopper-Pearson over pairs} & fails on a crossed design & AgentDojo, 28 pipelines, schemes (a); (b); (c) & 0.05 & \textbf{0.042--0.268; 0.060--0.340; 0.195--0.350} & 27 / 0 / 1; 28 / 0 / 0; 28 / 0 / 0 & 0.282 (26); 0.355 (27); 0.365 (28) \\
25 & \textbf{cluster bootstrap-t by user task, when injection tasks are sampled} & fails & AgentDojo, schemes (b); (c) & 0.05 & \textbf{up to 0.258; 0.275} & 21 / 0 / 7; 24 / 0 / 4 & 0.272 (20); 0.289 (24) \\
26 & \textbf{larger of the two clustered bounds} & holds under (a) only & AgentDojo, schemes (a); (b); (c) & 0.05 & at most 0.029; \textbf{up to 0.066; 0.089} & 0 / 0 / 28; 4 / 0 / 24; 16 / 3 / 9 & 0.034 (0); 0.074 (3); 0.098 (15) \\
27 & \textbf{pigeonhole bootstrap, basic limit} & fails & AgentDojo, schemes (a); (b); (c) & 0.05 & \textbf{up to 0.174; 0.178; 0.278} & 5 / 0 / 23; 14 / 1 / 13; 27 / 0 / 1 & 0.187 (5); 0.190 (10); 0.292 (26) \\
28 & \textbf{PPI on a stratified sheet read as an i.i.d. sample} & fails & three judge wordings, labels planted at three rates & 0.05 & \textbf{0--0.989} & 7 / 0 / 11 & 0.992 (7) \\
29 & \textbf{the judge's rate alone} & fails & the judge plasmodes (App.~\ref{app:judge}) & 0.05 & \textbf{0--1.000} & 30 / 0 / 12 & 1.000 (30) \\
30 & \textbf{multiway cluster variance with a t quantile} & fails & AgentDojo, schemes (a); (b); (c) & 0.05 & \textbf{up to 0.136; 0.175; 0.179} & 5 / 0 / 23; 8 / 3 / 17; 26 / 1 / 1 & 0.147 (5); 0.187 (7); 0.191 (25) \\
31 & \textbf{two clustered margins added in quadrature} & over in one cell under (c) & AgentDojo, schemes (a); (b); (c) & 0.05 & at most 0.015; 0.046; \textbf{0.058} & 0 / 0 / 28; 0 / 0 / 28; 1 / 1 / 26 & 0.020 (0); 0.053 (0); 0.066 (0) \\
\bottomrule
\end{tabular}

\caption{Miss rate against $\delta$, rows 19 to 31. Cells are counted as over / unresolved / at or under $\delta$; the next
column is the largest upper end of a 95\% interval for a cell's miss rate, with the cells still over after a
Bonferroni correction in brackets (for the exact binomial rows, the largest exact miss and the cells above $\delta$).}
\label{tab:fullB}
\end{table*}

\begin{table*}[p]
\centering
\scriptsize
\setlength{\tabcolsep}{3pt}
\begin{tabular}{@{}r p{2.8cm} p{2.0cm} p{3.9cm} l p{2.1cm} p{2.0cm} p{1.9cm}@{}}
\toprule
\textbf{\#} & \textbf{bound} & \textbf{kind} & \textbf{setting} & \textbf{$\delta$} & \textbf{miss, per cell} & \textbf{cells over / unresolved / at or under} & \textbf{largest 95\% upper limit of a cell's miss (cells over after correction)} \\
\midrule
32 & stratified Wald-t limit, 8 reference-rate strata & approximate & 5 mid-rate labels, real Granite-3.3-2B responses, $n_s$ 100--200, the draws of the StratPPI row & 0.05 / 0.10 & 0.007--0.036 / 0.027--0.080 & 0 / 0 / 28; 0 / 0 / 28 & 0.042 (0); 0.088 (0) \\
33 & the same limit, redrawn with replacement & approximate & the same 5 labels and sizes & 0.05 / 0.10 & 0.007--0.038 / 0.026--0.086 & 0 / 0 / 28; 0 / 0 / 28 & 0.040 (0); 0.088 (0) \\
34 & the same limit with the term for a sampled pool & approximate & a claim about the prompt source: the pool redrawn and the strata rebuilt, the same 5 labels and sizes & 0.05 / 0.10 & 0.007--0.050 / 0.031--0.095 & 0 / 1 / 27; 0 / 0 / 28 & 0.055 (0); 0.101 (0) \\
35 & \textbf{stratified Wald-t limit on a synthetic i.i.d. grid} & over in 5 of 18 cells at each $\delta$ & six strata profiles, $n_s$ 100--400, binomial strata & 0.05 / 0.10 & \textbf{0--0.111 / 0.086--0.147} & 5 / 5 / 8; 5 / 5 / 8 & 0.121 (4); 0.158 (4) \\
36 & \textbf{the Wilson-type limit on the same grid} & over in 1 and 4 of 18 cells & the same profiles and sizes & 0.05 / 0.10 & \textbf{0--0.057 / 0.086--0.135} & 1 / 2 / 15; 4 / 5 / 9 & 0.065 (0); 0.146 (2) \\
37 & stratified Wald-t limit inside the training loop, 8 strata & approximate & synthetic bandit, Lagrangian GRPO, 4 environments, the pool's own rate, 8,000 runs a cell & 0.05 / 0.10 & 0.029--0.052 / 0.072--0.100 & 0 / 1 / 3; 0 / 1 / 3 & 0.057 (0); 0.107 (0) \\
38 & the same limit with the sampled-pool term & over in 1 of 4 cells at each $\delta$ & the same runs, the population's rate & 0.05 / 0.10 & 0.047--0.057 / 0.093--0.107 & 1 / 1 / 2; 1 / 0 / 3 & 0.062 (1); 0.114 (0) \\
39 & stratified Wald-t limit, a costly label stratified on a cheap reference rate & over in 1 and 2 of 32 cells & XSTest human labels of five released models, 8 mid-rate labels, two stratifiers, $n_s$ 50--100, redrawn with replacement & 0.05 / 0.10 & 0--0.053 / 0.021--0.108 & 1 / 0 / 31; 2 / 0 / 30 & 0.055 (0); 0.111 (2) \\
40 & betting mixture on user-task rates; on the pair rate & exact (bounded) & AgentDojo, 28 pipelines, user tasks redrawn across the table (a) and within suites (a\_s); truth: the average over user tasks; the pair rate & 0.05 & at most 0.007; 0.005; 0.006; 0.003 & 0 / 0 / 28; 0 / 0 / 28; 0 / 0 / 28; 0 / 0 / 28 & 0.010 (0); 0.008 (0); 0.010 (0); 0.005 (0) \\
41 & cluster bootstrap-t, by user task, redrawn within suites & approximate & AgentDojo, 28 pipelines, scheme (a\_s) & 0.05 & 0.003--0.042 & 0 / 0 / 28 & 0.049 (0) \\
42 & \textbf{cluster bootstrap-t by user task, injection tasks redrawn within suites} & fails & AgentDojo, schemes (b\_s); (c\_s) & 0.05 & \textbf{up to 0.231; 0.249} & 11 / 1 / 16; 19 / 2 / 7 & 0.244 (11); 0.262 (18) \\
43 & two clustered margins in quadrature, redrawn within suites & approximate & AgentDojo, schemes (a\_s); (b\_s); (c\_s) & 0.05 & at most 0.016; 0.037; 0.040 & 0 / 0 / 28; 0 / 0 / 28; 0 / 0 / 28 & 0.020 (0); 0.044 (0); 0.046 (0) \\
\bottomrule
\end{tabular}

\caption{Miss rate against $\delta$, rows 32 to 43. Cells are counted as over / unresolved / at or under $\delta$; the next
column is the largest upper end of a 95\% interval for a cell's miss rate, with the cells still over after a
Bonferroni correction in brackets (for the exact binomial rows, the largest exact miss and the cells above $\delta$). Rows 32 to 36 are the \SWT{} and, on a
synthetic grid, it and the Wilson-type limit; 37 to 39 the same limit inside the training loop and on a costly human
label; 40 to 43 AgentDojo with an exact bound and with tasks redrawn within the benchmark's suites (schemes (a\_s), (b\_s) and (c\_s): as (a), (b) and (c), within each suite).}
\label{tab:fullC}
\end{table*}

\begin{table*}[p]
\centering
\scriptsize
\setlength{\tabcolsep}{3pt}
\begin{tabular}{@{}r p{2.8cm} p{2.0cm} p{3.9cm} l p{2.1cm} p{2.0cm} p{1.9cm}@{}}
\toprule
\textbf{\#} & \textbf{bound} & \textbf{kind} & \textbf{setting} & \textbf{$\delta$} & \textbf{miss, per cell} & \textbf{cells over / unresolved / at or under} & \textbf{largest 95\% upper limit of a cell's miss (cells over after correction)} \\
\midrule
44 & stratified Wald-t limit, a second model family, redrawn with replacement & approximate & Qwen2.5-1.5B-Instruct, refusal label (18--19\%), 8 reference-rate strata, steps 100 and 200, $n_s$ 100--200 & 0.05 / 0.10 & 0.033--0.046 / 0.077--0.092 & 0 / 0 / 4; 0 / 0 / 4 & 0.048 (0); 0.095 (0) \\
45 & the Wilson-type limit, the same & approximate & the same cells & 0.05 / 0.10 & 0.033--0.041 / 0.086--0.092 & 0 / 0 / 4; 0 / 0 / 4 & 0.043 (0); 0.095 (0) \\
46 & the stratified Wald-t limit with the term for a sampled pool, a second model family & over in 1 of 4 cells at $\delta$ 0.10 & a claim about the prompt source, the same label, steps and sizes & 0.05 / 0.10 & 0.036--0.051 / 0.085--0.111 & 0 / 1 / 3; 1 / 0 / 3 & 0.056 (0); 0.118 (1) \\
47 & stratified Clopper-Pearson: a stratified draw's pooled count plus the allocation's slack & exact (binary labels) & 116 cells: real responses (58), a costly label (40), a synthetic grid (18), drawn with replacement within strata, $n_s$ 50--400 & 0.05 / 0.10 & 0--0.048 / 0--0.093 & 0 / 0 / 116; 0 / 0 / 116 & 0.048 (0); 0.093 (0) \\
48 & stratified betting bound (union-intersection) & exact (bounded) & the same cells & 0.05 / 0.10 & 0--0.001 / 0--0.007 & 0 / 0 / 116; 0 / 0 / 116 & 0.002 (0); 0.010 (0) \\
49 & cluster bootstrap-t, by user task, against a fitted crossed model & approximate & AgentDojo, 28 pipelines, tables drawn from a fitted crossed model with injection tasks fixed (P\_a); truth: the model's rate & 0.05 & 0.001--0.044 & 0 / 0 / 28 & 0.051 (0) \\
50 & \textbf{the same, both kinds of task drawn from the model} & fails & the same, scheme (P\_c) & 0.05 & \textbf{up to 0.247} & 18 / 3 / 7 & 0.260 (17) \\
51 & betting mixture on user-task rates; on the pair rate, against a fitted crossed model & exact (bounded) & the same, scheme (P\_a); truth: the model's user-task average; its pair rate & 0.05 & at most 0.005; 0.003 & 0 / 0 / 28; 0 / 0 / 28 & 0.007 (0); 0.006 (0) \\
52 & pigeonhole bootstrap-t; two clustered margins in quadrature, against a fitted crossed model & approximate & the same, schemes (P\_a); (P\_c) for each & 0.05 & at most 0.011; 0.037; 0.016; 0.047 & 0 / 0 / 28; 0 / 0 / 28; 0 / 0 / 28; 0 / 0 / 28 & 0.014 (0); 0.043 (0); 0.021 (0); 0.054 (0) \\
53 & stratified Wald-t limit, a policy that training moved far, redrawn with replacement & approximate & Qwen2.5-1.5B-Instruct trained on the guard's refusal flag alone, refusal label (26\% at step 100, 58\% at step 200), the second family's pool and strata, $n_s$ 100--200 & 0.05 / 0.10 & 0.043--0.049 / 0.095--0.100 & 0 / 0 / 4; 0 / 0 / 4 & 0.051 (0); 0.103 (0) \\
54 & \textbf{the Wilson-type limit, the same} & over in all 4 cells at the 58\% rate & the same cells & 0.05 / 0.10 & \textbf{0.045--0.059 / 0.098--0.112} & 2 / 0 / 2; 2 / 1 / 1 & 0.062 (2); 0.115 (2) \\
55 & the stratified Wald-t limit with the term for a sampled pool, a policy that training moved far & approximate & a claim about the prompt source, the same label, steps and sizes & 0.05 / 0.10 & 0.040--0.048 / 0.085--0.094 & 0 / 0 / 4; 0 / 0 / 4 & 0.052 (0); 0.100 (0) \\
\bottomrule
\end{tabular}

\caption{Miss rate against $\delta$, rows 44 to 55. Cells are counted as over / unresolved / at or under $\delta$; the next
column is the largest upper end of a 95\% interval for a cell's miss rate, with the cells still over after a
Bonferroni correction in brackets (for the exact binomial rows, the largest exact miss and the cells above $\delta$). Rows 44 to 46 are the second model family after an update that kept its order of prompts, and 53 to 55 the same
family after a separate update, trained on the guard's flag alone, that reordered them; 47 and 48 exact stratified bounds; 49 to 52 AgentDojo against a fitted crossed model (schemes $P_a$, injection
tasks fixed, and $P_c$, both drawn), where only a bound that fails to compute counts as a miss.}
\label{tab:fullD}
\end{table*}

\paragraph{Abstention.} A bound that returns no limit, or a limit of 1, gives no certificate and cannot miss. A reader
who sees only the limits that were returned faces the conditional miss rate, the misses divided by the draws that
returned a limit. Computed afterwards from the stored shares, it was over in no cell of the bootstrap-t rows of
Table~\ref{tab:valid}: its largest value was 0.054 for PPI++, 0.056 for StratPPI and 0.051 for the bound clustered by
user task, because the cells that abstain most are those that miss least. The
bootstrap-t StratPPI limit returns none in up to 6\% of draws on the reference-rate strata (up to 8.6\% redrawn with
replacement) and in up to 52\% on the judge-logit strata at a 1.3\% rate; PPI++ with a bootstrap-t limit in up to 55\% of
draws where a sample of 225 holds about three positives; the clustered AgentDojo limits in up to 82\% of tables (the pigeonhole
bootstrap-t; the cluster bootstrap-t by user task in up to 15\% on these draws), for the pipelines with the fewest
successes. Counted as misses, the rows that hold would be over in up to 7 of 28 cells (row 12
at $\delta=0.05$), in 10 of 168 for row 9 and for 2 to 7 of 28 pipelines in the AgentDojo rows. Counted this way even an
exact bound can exceed its level: Clopper-Pearson returns 1 when every label is a positive, which at a label with rate 0.9535 and a
safety set of 100 happens in 0.9\% of draws and takes it to 0.052 at $\delta=0.05$. The count measures how much of a
record is abstention, not validity.

\paragraph{The Wilson limit, exactly.} For an i.i.d.\ sample the Wilson limit's miss probability, computed by
enumeration, jumps just above each value the limit can take and falls until the next; averaged over rates it is under
$\delta$ below one half and over it above. At $n=100$ and $\delta=0.05$ it averages 0.042 over rates of 5--20\% and 0.056
over 80--95\%. Redrawn with replacement, the \SWL{} is over its level in 7 of 28 mid-rate cells at $\delta=0.05$, all on
the two labels at rates of 65\% and above, and in none of the 16 at rates of 9--18\%; the pooled Wilson bound is over in
6 of 28. Rare labels take Clopper-Pearson or a betting bound.

\section{Selection and the Training Loop}
\label{app:pipeline}

The resampling studies hold the candidate fixed. Two checks test the gate with selection in it. With the training code of
our language-model runs (GRPO with a Lagrangian on the judged rate, a predicted safety test choosing the checkpoint)
applied to a synthetic bandit at four levels of between-prompt heterogeneity, 8 strata and 8,000 runs a cell, the \SWT{}
for the pool's own rate was over its level in none of its 8 cells (2 unresolved; largest miss 0.052 at $\delta=0.05$
and 0.100 at 0.10). With the sampled-pool term, for the population's rate, it was over in the least heterogeneous
environment at both deltas (0.057 and 0.107) and in no other cell; the Wilson-type limit with the term and the pooled
bound on a random split were over in none. With 4 strata and 500 runs a cell, the \SWT{} was over in one of four cells
(0.132 at $\delta=0.1$). In a classifier setting with a constraint on a difference of true-positive rates,
rerandomised and stratified candidate/safety splits, across eleven split rules, 2,000 seeds per arm and candidates built
to overfit the balancing statistics, returned a passing, truly violating candidate in at most 1 of 2,000 runs under
Clopper-Pearson; a tight Wald test on the same runs, sensitive enough to show a leak, fell below the truth in at most
0.055 of runs at $\delta=0.05$.

\section{Strata from the Pre-Update Model in Detail}
\label{app:strata}

\paragraph{First pools.} With $k=8$ and $H=8$ on Granite-3.3-2B's responses, $\mathrm{ESS}_{\text{est}}$ at $\delta=0.10$
was 2.4 for over-refusal (rate 17\%, $\mathrm{ICC}_{\text{ref}}$ 0.72), 5.1--5.3 for refusal of plainly harmful requests
(66\%, 0.86) and 1.4 for non-refusal of encoded requests (9\%, 0.50); placebo strata cut on a permuted reference rate gave
1.00 to 1.04 with 4 strata on every label but the one at 93\%. When the Lagrangian pushed the stratified label itself,
the gain was 2.1--2.2, with no compression of the between-prompt spread and the order of prompts kept ($\omega$ 0.92),
in a run whose rate ended within 0.4 to 3.4 points of the reference's. The rule registered in advance collapsed to one to
three strata, because most prompts have an 8-sample reference rate of exactly 0 or 1; equal rank strata with random ties
were adopted on the same pools, and the truth used for coverage, which carried a fixed offset shared by every draw,
became the mean of the half of the responses the draws come from; these gains were thus measured on the data that chose
the rule (Appendix~\ref{app:record}).

\paragraph{The pre-flight} ranks the gain (Spearman 0.84 between predicted and realised $\mathrm{ESS}_{\text{est}}$ over
432 cells at $\delta=0.10$) and over-predicts its size. On new prompts it was within 20\% at all three mid rates with the
updated policy's two quantities measured on responses the draws do not use; from the reference samples alone it
predicted 2.58, 4.12 and 2.22, and the realised gains were 27\% to 41\% lower.

\paragraph{Against StratPPI.} On the same strata and draws, StratPPI as published is shorter partly because its normal
limit runs over its level (Table~\ref{tab:stratppi}); with a bootstrap-t limit and proportional allocation its estimator
holds, with a median ESS of 2.10 against 2.16 for the Wilson-type limit. At a safety set of 100 with 8 strata the
within-stratum regression on the reference rate adds nothing once the limit is valid; at 200, or with 4 strata, it adds
0.4 to 0.8. For a claim about the prompt source, with the pool redrawn and the strata rebuilt, limits without the
sampled-pool term fail: the Wilson-type limit is over in 17 of 28 cells and the bootstrap-t StratPPI limit in 21.

\begin{table*}[t]
\centering
\scriptsize
\setlength{\tabcolsep}{3pt}
\begin{tabular}{@{}p{4.2cm} r p{2.4cm} p{2.6cm} p{3.4cm}@{}}
\toprule
\textbf{label (rate)} & \textbf{$n_s$} & \textbf{the Wilson-type limit} & \textbf{StratPPI as published} & \textbf{StratPPI estimator, bootstrap-t limit} \\
\midrule
over-refusal (0.16) & 100 & 0.023; 2.39 & \textbf{0.066}; 4.11 & 0.042; 2.58 \\
over-refusal (0.16) & 200 & 0.024; 2.30 & 0.052; 3.46 & 0.042; 3.10 \\
over-refusal, pushed by the Lagrangian (0.18) & 100 & 0.024; 2.19 & 0.0558; 3.24 & 0.036; 2.43 \\
over-refusal, pushed (0.18) & 200 & 0.023; 2.13 & 0.045; 2.85 & 0.037; 2.49 \\
non-refusal, encoded prompts (0.09) & 100 & 0.034; 1.38 & \textbf{0.086}; 2.59 & 0.045; 1.43 \\
non-refusal, encoded prompts (0.09) & 200 & 0.027; 1.36 & \textbf{0.063}; 2.11 & 0.041; 1.59 \\
refusal, harmful prompts (0.65) & 100 & 0.047; 4.74 & \textbf{0.0564}; 6.30 & 0.033; 1.77 \\
refusal, harmful prompts (0.65) & 200 & 0.040; 5.19 & 0.045; 6.29 & 0.034; 5.35 \\
refusal, encoded prompts (0.93) & 100 & 0.054; 0.96 & 0.032; 0.78 & 0.031; 0.83 \\
refusal, encoded prompts (0.93) & 200 & 0.045; 1.12 & 0.026; 0.92 & 0.033; 1.03 \\
\bottomrule
\end{tabular}

\caption{Reference-rate strata at $\delta=0.05$: largest miss over checkpoints; ESS against a random split with a
pooled Wilson bound. The rates are those of these draws; the first-pools paragraph quotes an earlier measurement of the same labels, from the estimate at $\delta=0.10$, whose rates differ by a point. Bold: more than two Monte Carlo standard errors over $\delta$ (above 0.0562 at 5,000 draws a cell;
the two cells nearest it carry a fourth digit).}
\label{tab:stratppi}
\end{table*}

\paragraph{The registered confirmation.} On three new pools (refusal rates 12\%, 82\%, 98\%), six of eight predictions
were kept. With a large pool the Wilson-type limit was over its level in none of 4 cells at 12\% and in all 4 at 82\%;
the bootstrap-t StratPPI limit was over in 1 of 8 mid-rate cells (0.105 at $\delta=0.10$, at 82\%) and in one at 98\%
(0.0524 at $\delta=0.05$); the \SWT{} in none of 8. For the source, the Wilson-type limit with the term was over at 82\%
with a safety set of 100 (0.058), where the \SWT{} with the term was over in none of 8. On the mixed pool at 46\%, nine of
ten predictions were kept and one could not be tested; it is weaker evidence, because its rate was aimed at and its
predictions were written after a dry run on a re-mix of the first pools' responses, and the reference rate largely
sorts its two sources into different strata. The second family's update, meant to reorder prompts, did not
(correlation 0.84; three of six predictions kept). A separate update of the same model, trained on the guard's flag
alone, did reorder them
(0.46): the \SWT{} stayed valid and gained 0.97 and 1.08, the Wilson-type limit was over in all 4 cells at 58\%, and four of
five predictions were kept; the one refuted is the check of the push's direction (refusal rose on the categories it was
meant to lower). By its registration, a failed direction check meant the policy had not moved; the rate nearly tripled
and the order of prompts moved further than in any other run, so we report it as moved, against that reading.

\paragraph{A costly label.} With XSTest's human labels on five released models' completions for its 250 safe prompts,
strata from the guard's rate over 8 Granite-3.3-2B reference responses, safety set 100 drawn with replacement, the
\SWT{}'s gain over eight mid-rate labels was 0.97 to 1.34 (median 1.18); it was over its level in 3 of 64 cell-and-level pairs by at most
0.008 (the skew-corrected limit: 0.93 to 1.21, median 1.02, over in none). The Wilson-type limit gained 1.10 to 1.36 but was over in 8 of 32 cells at $\delta=0.05$.

\paragraph{Exact stratified bounds and the chance of certifying.} Over 116 cells drawn with replacement within strata,
Clopper-Pearson on the stratified draw's pooled count plus the allocation's slack, and a stratified betting bound
\citep{spertus2022sweeter}, kept their level everywhere. At a threshold where Clopper-Pearson on a random draw passes with
probability 0.8, the stratified Clopper-Pearson bound's power ESS on the 36 real-response cells at rates of 5\% to 95\%
was 1.11 (median; 0.27 to 1.83) and on the 32 costly-label cells 0.88; the betting bound's was 0.05; the \SWT{}'s 1.88.

\paragraph{The rule in one place.} For 8 equal rank strata of an 8-sample reference rate: a label that may be rare (under
about 5\%) or near one (above about 95\%) takes no strata and Clopper-Pearson or a betting bound; a mid-rate label (about 5\%
to 80\%) takes the skew-corrected \SWT{} (next paragraph; gains seen at 12\% to 46\%, none on human labels), with the sampled-pool term for a claim about the prompt source; from about 80\% to 95\% a random
draw serves as well; where a finite-sample guarantee is required, Clopper-Pearson on a random draw; where the update may
reorder the prompts, the same limits, expecting little gain.

\paragraph{Why the \SWT{} misses, and the skew correction.}\label{app:skew} The limit is centred on the estimate, and at
rates under one half the estimate is right-skewed. The one-term Edgeworth expansion of the studentised estimate
\citep{hall1992,johnson1978} puts the miss at $\delta+\gamma(2z^2+1)\phi(z)/6$, with $\gamma$ the estimate's skewness.
Adding $\max(0,\hat\kappa_3/\hat V)(2t^2+1)/6$ to the limit, with
$\hat\kappa_3=\sum_hW_h^3\tilde p_h(1-\tilde p_h)(1-2\tilde p_h)/n_h^2$ and $\hat V$ the limit's own variance, cancels
that term to first order and never lowers the limit. We wrote this correction from the expansion, with no tuned constant, after seeing
where the limit failed. On the 116 cells of the exact-bound comparison above, rerun at 40,000 draws a cell and both
levels, the uncorrected limit was over in 19 of 232 cell-and-level pairs and the corrected one in 5, all on the
synthetic grid (where, at this number of draws, the uncorrected limit was over in 8 of the 18 cells at each level); on the real-response and costly-label cells the corrected limit was over in none of 196. We then
registered six predictions, pushed before the run, on a new synthetic grid of 648 cells: four shapes of the strata
rates, 2, 4 and 8 binomial strata, rates of 2\% to 90\%, equal and unequal shares, safety sets of 100 to 400. Four were
kept and two refuted. At rates of 5\% to 50\% the uncorrected limit was over in 339 of 720 pairs (47\%), by up to 0.069,
and the corrected one in 76 (11\%), by up to 0.023, at a median 0.90 and 0.91 of the uncorrected limit's gain at the two
levels; under 5\% the corrected limit was over in 8 of 144 pairs, all with 2 strata at $\delta=0.10$, and Clopper-Pearson remains the rule there. Refuted: the corrected limit's largest miss at
$\delta=0.10$ and rates of 5\% to 50\% was 0.123, against a predicted 0.12 at most, and above 50\% it was over in 12 of
432 pairs (2.8\% against a predicted 2\%), with a gain as low as 0.88 of the uncorrected limit's where we predicted a loss of at most 1\%. Skew does not explain the uncorrected limit's 33 pairs over above 50\%. By number of strata, a breakdown made after the run, at rates of 5\% to 50\%: with 2 strata the corrected limit
was over in 57 of 240 pairs, with 4 in 12 of 240 and with 8 in 7 of 240 (by at most 0.007), where the uncorrected limit
was over in 79 of 240 (by at most 0.023). Cells are counted by their nominal rate, a correction to the scoring code after the run that changed no verdict. The rule registered with the predictions reads this as a partial repair and keeps the uncorrected limit. We
override that rule, a decision taken after the result and so not a registered one: the corrected limit is never below
the uncorrected one and fails far less often. We then recomputed, without registration, the cells behind Table~\ref{tab:strata} for the
corrected limit, on the draws of the registered runs for the pool's own rate and on new draws of the same design for the
prompt source: over 168 cell-and-level pairs the corrected limit was over its level in none, the
uncorrected one in 5. With the sampled-pool term the correction divides the second phase's third moment by the total
variance. Table~\ref{tab:strata} gives the corrected limit's gains; every other gain and miss rate of the \SWT{} in this
appendix, Tables~\ref{tab:fullC} and~\ref{tab:fullD} included, is the uncorrected limit's. On the human labels the
corrected limit's gain at a safety set of 100 was 0.93 to 1.21 (median 1.02). We did
not establish why the uncorrected limit kept its level on the real-response cells of that comparison and not on binomial
strata of the same sizes.

\section{One Guarantee per Attack Family in Detail}
\label{app:family}

\paragraph{Set-up.} A simulated label picks a prompt uniformly within its family and a 0/1 label with that prompt's
measured rate over its 24 guard-judged responses, so each family's true rate is the mean of its prompts' rates. The
gated label counts a response to an encoded prompt only if it is also similar to the request (similarity at least 0.8)
and flagged unsafe; the ungated label is the guard's flag alone. A family is certified safe if its upper bound is at most
$\tau$, unsafe if its lower bound exceeds $\tau$, and indifferent if both bounds lie within $\tau\pm0.4\tau$, allowed only
for a family whose true rate lies there. Fixed-size designs use Clopper-Pearson at $\delta/7$ on a grid of 100 to 20,000
labels; the sequential designs use 10 labels per family to start, a beta-binomial mixture bound with a uniform prior at
$\delta/7$ (valid under adaptive sampling and stopping by Ville's inequality), and a cap of 20,000 labels. Both
simulations were registered before they were run; every registered prediction held. They were run after the tally of
Appendix~\ref{app:record} was closed and are not among its 72 predictions.

\paragraph{All cells.} Labels to resolve all families in 80\% of runs (uniform / equal / index rule): gated, $\tau=0.02$:
3,000 / 10,000 / 5,809; $\tau=0.05$: 7,000 / 5,000 / 2,441; $\tau=0.10$: 10,000 / 7,000 / 2,767; ungated, $\tau=0.02$:
over 20,000 / over 20,000 / 14,643; $\tau=0.05$: over 20,000 / 20,000 / 18,765; $\tau=0.10$: 10,000 / 5,000 / 3,563. On a finer grid (steps of 250, not registered) the registered cell's fixed-size counts are about 6,500 and 4,000. The
finish-first rule needed the same totals as the index rule within 1.5\% in every cell and resolved at least as many
families as uniform at 600 labels in every cell but the ungated $\tau=0.05$ one (1.8 against 2.1). Equal labels per family with Clopper-Pearson, run
afterwards and not registered, resolved 0.1, 5.6, 5.9, 6.3 and 6.6 families at 300, 600, 1,000, 1,500 and 2,500 labels in
the registered cell, against 1.0, 2.8, 4.6, 6.0 and 6.8 for the finish-first rule. The index rule's
numbers in Table~\ref{tab:family} come from the second simulation (2,444 there, 2,441 in the first).

\paragraph{A learned sampler, tried first.} We first trained a second adapter of the same model to propose harmful
responses, as an importance-sampling proposal for the breach rate. Its forward divergence from the model's own harmful
responses had median 187.2 nats on 41 prompts, against 2.23 for the model itself: it wrote short harmful outputs the
model rarely writes, and a supervised start was as far, so the training objective was not the cause. Moved to prompt space (a scorer of prompts mixed with
the prompt distribution), its bound kept its level (coverage 0.99 to 1.00) but was 2.4 times wider than Clopper-Pearson
on a random draw at 300 labels in its most favourable registered cell (0.0308 against 0.0126; up to 6 times at a
smaller mixture weight), no narrower than with uniform prompt scores, and its scores ranked new prompts at Spearman
0.301. That study used another model (Qwen3-1.7B), another judge (WildGuard) and a pool of adversarial and benign
prompts, and it certified the average, so it is not a like-for-like comparison with the families here.

\section{Judges in Detail}
\label{app:judge}

\begin{table}[t]
\centering
\small
\begin{tabular}{@{}p{3.4cm}p{3.6cm}@{}}
\toprule
data & route \\
\midrule
a property code can verify & compute it on every response; Clopper-Pearson (exact) \\
gold labels, fewer than 10 in the rarer class & Clopper-Pearson on the labels; judge unused (exact) \\
gold labels, 10 or more in each class & PPI++ on the judge's logit, bootstrap-t limit (approximate) \\
gold labels from a stratified sheet & design-weighted labels, skew-corrected \SWT{} (approximate; the rule was settled with the Wilson-type limit, replaced after it ran over its level) \\
a calibration measured on another population & refused \\
\bottomrule
\end{tabular}
\caption{The routing rule, fixed before a constraint's labels are read. It was settled on the resampling results (20
cells) and has had no confirmation run. On the one stratified sheet design tested the
Wilson-type limit was over its level in 1 of 18 cells (0.060), and it can be over once the pool is large
(Section~\ref{sec:valid}); the skew-corrected \SWT{} is the safer choice
there, and an exact stratum-wise Clopper-Pearson construction (Section~\ref{sec:human}) where a guarantee is required.}
\label{tab:routing}
\end{table}

\paragraph{What a judge is worth.} For the compiled refusal judge against the guard's label (225 labels, 2,000 judged
draws), ESS against Clopper-Pearson on the labels is 1.11 for the 0/1 verdict, 1.15 for the probability, 1.61 for the
logit and 1.86 for a cross-fitted Platt score \citep{platt1999}. The logit carries more than the verdict.

\paragraph{Labels drawn within strata of the judge.} With labels allocated in proportion to stratum size, StratPPI's
estimator with a bootstrap-t limit was over its level in none of 28 cells (5 unresolved, largest miss 0.056; 5 would be
over if its returns of no limit counted as misses, 4 of them at the 1.3\% rate), with a median ESS of 2.54 at a 20\% rate
and 1.87 at 5\% with 10 strata, against 1.67 and 1.39 for unstratified PPI++ with the same limit. The 0.90 cell of
Table~\ref{tab:valid} is the rubric judge at a 5\% rate, 1,000 labels, 5 strata and the heuristic allocation: four
strata with judge probability under 0.001 hold 19\% of the positives at rates of 1.0--1.5\%, and their eight labels show
no positive in 91\% of draws.

\paragraph{The carried calibration, by seed.} Across the targeted update, recall differed (Fisher's exact test, $p<0.05$, unadjusted) in 2, 2 and 1 of six
wordings and the carried bound was over its level for 3, 3 and 2, because it also fails where a fall in recall is not significant (seed 0, one wording: 0.04 to 0, $p=0.12$, miss 0.23) and, in seed 2, where the false-alarm rate differed (seed 0: misses 0.08, 0.23 and 0.80; the largest in
seeds 1 and 2 are 0.399 and 0.625); at step 100 it was over for 1, 1 and 4. Seeds 1 and 2 returned \NSF{} at their gate
(upper bounds 0.2317 and 0.2064 against a threshold of 0.2058; seed 0's step 175 had 0.1979 and was returned), so the
step-200 checkpoints analysed are not policies the gate returned; the question is what a carried calibration says about
an updated policy, which does not depend on that. We wrote down beforehand, in a note with no independent timestamp, that
a failure at step 200 would mean the earlier success was luck and no calibration carries across training; reading the
side-effect run as a control came after the result. WildGuard, as judge with the guard as gold, had false alarms of 1.0\%
before and 1.0\% to 1.7\% after; against the annotator's strict refusals after the update, WildGuard and the guard agree
equally well (design-weighted accuracy 0.930 and 0.929), and the compiled judge flags none of 18. WildGuard's own
calibration did not carry from the harmful pool to the benign one (85\% misses).

\paragraph{The gold label is the guard's.} The guard is not the truth on these responses: of 18 guard-flagged long refusals of the updated policy, the annotator called 12 refusals and 6 answers, so part of the fall in the judge's recall may be a change in what the guard flags. The guard's precision against the annotator is 0.59 before and 0.60 after, and with WildGuard as gold the judge's recall still fell, from 0.23 to 0.03, 0.08 and 0.05; but WildGuard flags 95\% of the guard's long refusals, so only the 18 human-labelled refusals are independent of both.

\paragraph{The mechanism we suspect.} By the guard and a surface pattern (a refusal opener and more than 40 words), on
the 500-response sets that both judges scored, 29 of the pre-update model's 100 guard flags are long refusals, against
83 of 93, 83 of 108 and 85 of 97 after the update in seeds 0, 1 and 2. Of the 93 guard-flagged refusals of seed 0 after
the update, the rubric judge called 90 answers. Of 18
such responses, the annotator labelled 12 refusals and 6 answers. WildGuard flagged 95\% of the guard-flagged responses
with the pattern and 92\% of those without. The rubric judge flagged 10\% of the pre-update model's guard-flagged
responses with the pattern and 25\% of those without (canonical wording); after the update, about 14 of 251 with it, against 3
of 29 before (no detectable change), and about 2 of 47 without it, against 18 of 71 before, so training changed both how many responses
have the pattern and how the judge reads the plain ones. The pattern is not a label, and the
two runs also differ in reward and set-up.

\paragraph{Cutting responses to their first sentence.} A check registered before any cut response was scored tests the
body directly. Every response of the pre-update set and of the three seeds' step-200 sets (500 each) was cut to its first
sentence (the text up to the first full stop, exclamation mark or question mark followed by white space or the end; failing that, the first line) and scored again by
the same judge under the same six wordings; gold stayed the guard's flag on the full response. All five predictions
held. Among the 251 guard-flagged responses with the pattern after the update, the canonical wording flagged 14 in full
and 156 once cut, and lost none (exact McNemar $p<10^{-40}$). Its recall after the update rose from 0.03, 0.07 and 0.05
to 0.57, 0.56 and 0.56, and before the update from 0.21 to 0.54, so the body hides refusals from the judge on both
sides of the update, consistent with the calibration failing because the share of refusals with a body changed. With both sides cut,
the carried bound was over its level for 1 of six wordings in each seed, against 3, 2 and 2 with full responses on
these draws (seed 1's second wording missed in 0.057 of draws, at the line between unresolved and over, and was over on the draws
of Section~\ref{sec:carry}). The cut bound holds mostly by being loose: its median is 0.36 to 0.41 under the canonical
wording and 0.60 to 1 under the others, against rates of 0.19 to 0.22; with only the updated policy's responses cut it
returned no certificate at the median under any wording, so the prediction about that arm says little. Cutting also
raised false alarms on the guard's non-refusals, from at most 0.04 to between 0.20 and 0.57 under four wordings (to at
most 0.03 under the other two). The rule is crude (it stops at an abbreviation) and was not tuned. \paragraph{A rubric that names the long refusal.} All six compiled rubrics exclude a response that answers or attempts to answer, even in part (one names ``relevant
information''), so the judge's definition is narrower than the guard's. A second registered check asked whether
that explains the failure. One new sentence (``The response refuses the request or declines to answer it, even if the
refusal is followed by an explanation, advice or other related information'') was compiled by the same model with the
same prompt and scored on the full responses of the same four sets. The compiled rubric partly undoes the sentence: it also says the property does not hold if the response ``includes the refusal as part of a broader answer''. Four of five predictions held. Its recall on the
guard's refusals was 0.76 before the update (76 of 100) and 0.78, 0.70 and 0.62 after (73 of 93, 76 of 108 and 60 of 97),
with false alarms of at most 0.01 throughout, and it flagged 67 of 83, 63 of 83 and 57 of 85 of the long refusals. The
carried bound missed in 0.000, 0.007 and 0.072 of draws (over above 0.057): over its level in seed 2, where recall fell (Fisher $p=0.045$),
against our prediction of no seed; its medians were 0.28, 0.29 and 0.23 against rates of 0.19, 0.22 and 0.19. The sentence was written with hindsight, on these sets, knowing which responses the
earlier judge missed, and is one wording without paraphrases; even so its calibration did not carry in one seed of three.

These five predictions, the five of this second check,
and the seven of the two AgentDojo checks of Appendix~\ref{app:bench}, were registered after the tally of
Appendix~\ref{app:record} was closed and are not among its 72.

\section{Benchmarks in Detail}
\label{app:bench}

\paragraph{AgentDojo.} The published runs hold 19,380 episodes. Table~\ref{tab:agentdojo} gives six pipelines. The
bound clustered by user task returns no limit when the $\delta$-quantile of its bootstrap is not finite: on the fresh
draws of a registered run, in 9\% of the tables of the pipeline that certifies and in 16\% for the two Meta-SecAlign-70B
pipelines. Redrawing injection tasks within the benchmark's suites, the bound is over for 11 of 28 pipelines. Against
a fitted crossed logit-normal model with independent task effects, the same bound is over for none with injection
tasks fixed and for 18 of 28 with both kinds of task drawn; the multiway limit for 13 and the pigeonhole bootstrap-t for
none. Under the model, the pipeline that certifies on its published table certifies in 99\% of tables with the injection
tasks fixed under the bound clustered by user task, and in 35\% under the betting bound on user-task rates. The
certificate's bound is sensitive to the bootstrap's size when a handful of clusters hold every success: for the
defended model, at 4,000 draws it ranged from 0.058 to 0.107 over 30 seeds; the validity simulations use 4,000 draws inside each of 4,000 resampled tables, and the reported limits 200,000.

\begin{table*}[t]
\centering
\scriptsize
\setlength{\tabcolsep}{3pt}
\begin{tabular}{@{}p{3.0cm} r r r r r l r p{1.6cm}@{}}
\toprule
\textbf{pipeline} & \textbf{pairs} & \textbf{rate} & \textbf{per-pair bound} & \textbf{clustered by user task} & \textbf{clustered by injection task} & \textbf{``at most 5\%''} & \textbf{exact, user-task average} & \textbf{both kinds of task sampled} \\
\midrule
claude-3-5-sonnet-20241022 & 629 & 0.011 & 0.021 & 0.022 & 0.035 & pass & 0.048 & no limit \\
Meta-SecAlign-70B & 949 & 0.022 & 0.032 & 0.062 & 0.035 & \NSF{} & 0.092 & 0.091 \\
command-r & 629 & 0.033 & 0.048 & 0.063 & 0.054 & \NSF{} & 0.071 & 0.078 \\
claude-3-7-sonnet-20250219 & 949 & 0.050 & 0.063 & 0.070 & 0.104 & \NSF{} & 0.119 & 0.107 \\
gpt-4o + tool filter & 629 & 0.068 & 0.087 & 0.092 & 0.106 & \NSF{} & 0.107 & 0.115 \\
gpt-4o-2024-05-13 & 629 & 0.477 & 0.511 & 0.537 & 0.588 & \NSF{} & 0.593 & 0.604 \\
\bottomrule
\end{tabular}

\caption{Targeted attack success on AgentDojo's published runs, upper bounds at $\delta=0.05$. The certificate is the
bound clustered by user task with injection tasks taken as fixed; the verdict column is for it. The exact column is the
betting bound on the user tasks' success rates, which bounds their average; the last column is the pigeonhole
bootstrap-t, which also takes injection tasks as sampled. Bootstrap bounds from 200,000 draws.}
\label{tab:agentdojo}
\end{table*}

\paragraph{Two registered checks of the AgentDojo certificate.} \emph{A closed form.} The cluster t-limit by user
task with $m-1$ degrees of freedom, and the same limit with the one-term skew correction of
Appendix~\ref{app:skew} computed from the user tasks' third moment, need no bootstrap and returned no limit in at most
0.6\% of resamples. With user tasks resampled and injection tasks fixed (20,000 resamples per pipeline), the plain
limit was over its level for 24 of 28 pipelines (misses up to 0.139) and the corrected one for 11 (up to 0.083), with 9
more unresolved; we had predicted at least 8 and at most 4, with misses of at most 0.08. On the real tables the corrected limit is
0.021 for the pipeline that certifies (bootstrap-t 0.022, as predicted) but also passes both
Meta-SecAlign-70B pipelines at 5\% (0.049 and 0.047), against our prediction of one pass; it is over its level for
all three pipelines it passes (misses 0.080, 0.071 and 0.083), and the bootstrap-t refuses the two (0.061 and 0.057). The closed form does not replace the
bootstrap-t here. \emph{Leaving one injection task out.} For each pipeline, suite and injection task, we computed the
bootstrap-t by user task on the suite's other injection tasks and compared it with the left-out task's own success rate
over the same user tasks. The left-out rate was above the limit in 163 of the 638 cases with a limit (26\%; 161 of 799
had none), in 25 of 28 pipelines (predicted: more than 15\%, at least 20). For the pipeline that certifies, two tasks
succeed in 2 of 16 and 3 of 21 pairs (counts we had seen before registering; neither case returned a limit, and none of its 11 with a limit was exceeded). A limit on a mean over tasks is not a
limit on one task, so these are not misses; they show how little the certificate says about a task it has not seen.

\paragraph{A robot benchmark.} RoboDojo-RC Tier 1 \citep{machcha2026robodojo} publishes 120 trials per model, 20 on each
of six tasks the report selected; the harness's overheat termination is the label. With the six tasks taken as fixed,
the pooled Clopper-Pearson limit is valid for their mean \citep{hoeffding1956successes}. Table~\ref{tab:robot} shows what
120 trials resolve: a 2.5\% rate when nothing is observed, and not 5\% after two events; a 1\% claim with nothing
observed needs 299 trials. Taken as a sample of tasks, six tasks certify nothing useful. In an exploratory analysis, the
peak joint effort the rig reports predicts a stop (AUC 0.75 to 0.83 by model), partly mechanically, since the stop is an
overheat.

\begin{table*}[t]
\centering
\scriptsize
\setlength{\tabcolsep}{2pt}
\resizebox{\ifdim\width>\textwidth\textwidth\else\width\fi}{!}{\begin{tabular}{@{}l r r r r r p{2.6cm} r@{}}
\toprule
\textbf{model} & \textbf{n} & \textbf{stops} & \textbf{rate} & \textbf{upper, $\delta$ 0.05} & \textbf{upper, $\delta$ 0.10} & \textbf{``at most 5\%'' at $\delta$ 0.05} & \textbf{tasks as a sample} \\
\midrule
Claude Opus 5 & 120 & 10 & 0.083 & 0.137 & 0.125 & \NSF{} & 0.480 \\
Claude Opus 5.5 & 120 & 2 & 0.017 & 0.052 & 0.044 & \NSF{} (passes at $\delta$ 0.10) & 0.441 \\
GPT-6 Astra & 120 & 2 & 0.017 & 0.052 & 0.044 & \NSF{} (passes at $\delta$ 0.10) & 0.442 \\
\bottomrule
\end{tabular}
}
\caption{Safety stops on a robot benchmark, upper bounds; the last column takes the tasks as a sample (betting bound
on six task rates).}
\label{tab:robot}
\end{table*}

\section{Human Labels and the Cost of the Relative Certificate}
\label{app:human}

\begin{table*}[t]
\centering
\scriptsize
\setlength{\tabcolsep}{3pt}
\begin{tabular}{@{}p{4.0cm} p{2.6cm} r p{2.6cm} p{2.6cm}@{}}
\toprule
\textbf{model} & \textbf{humans: full / partial refusal} & \textbf{guard flags} & \textbf{strict: recall, false alarms} & \textbf{broad: recall, false alarms} \\
\midrule
GPT-4 & 0.064 / 0.020 & 0.088 & 1.00, 0.03 & 0.90, 0.01 \\
Llama 2, original system prompt & 0.380 / 0.216 & 0.636 & 0.96, 0.44 & 0.94, 0.19 \\
Llama 2, no system prompt & 0.140 / 0.156 & 0.304 & 0.94, 0.20 & 0.85, 0.07 \\
Mistral instruct & 0.008 / 0.008 & 0.036 & 0.50, 0.03 & 0.50, 0.03 \\
Mistral with a guardrail prompt & 0.096 / 0.092 & 0.168 & 0.75, 0.11 & 0.70, 0.04 \\
pooled (1,250) & 0.138 / 0.098 & 0.246 & 0.92 (0.88--0.96), 0.14 & 0.87 (0.82--0.91), 0.05 \\
\bottomrule
\end{tabular}

\caption{The guard's refusal flag against XSTest's human labels on safe prompts: recall and false-alarm rate against full
refusal (strict) and against full or partial refusal (broad). Pooled row: 95\% bootstrap intervals over prompts.}
\label{tab:xstest}
\end{table*}

The guard marks 92\% of human full refusals (159 of 172), 80\% of partial refusals and 5\% of full compliance; its
agreement with the strict human label (Cohen's kappa 0.59) is below two humans' (0.90). The squared correlation of the
strict label with the guard's logit is 0.46, roughly a twofold gain in labels as a predictor.

\begin{table*}[t]
\centering
\scriptsize
\begin{tabular}{@{}p{5.6cm} p{3.0cm} p{4.6cm}@{}}
\toprule
 & \textbf{reference model} & \textbf{policy trained under the constraint (step 200)} \\
\midrule
the guard's refusal flag, share of the pool & 17.9\% & 17.4\% \\
human label: refuses & 10.7\% [4.7, 21.5] & 10.6\% [5.2, 21.2] \\
human label: refuse-then-answer & 3.4\% [0.4, 15.5] & 0.3\% [0.0, 11.0] \\
guard's false-alarm rate against ``refuses'' & 8.1\% & 7.7\% \\
guard's precision & 0.59 & 0.60 \\
guard's recall: estimate, conservative lower limit & 0.99, 0.44 & 0.98, 0.45 \\
$\rho^2$ of the human label with the guard's logit & 0.61 & 0.60 \\
\bottomrule
\end{tabular}

\caption{One annotator's labels of 220 responses (conservative 95\% intervals).}
\label{tab:human}
\end{table*}

\paragraph{Our policies.} The annotator labelled a stratified sample of 220 responses to benign prompts of the pre-update
model and the updated policy (Table~\ref{tab:human}), under a guideline that follows XSTest's classes; the responses were
cut at 128 tokens, so a refusal followed late by an answer would be labelled a refusal, and the post-update policy
writes longer refusals; the table's refuse-then-answer share may be undercounted for it. No difference in strict refusal was detected ($-0.05$ points, prompt-clustered standard error 3.9).
The certificate of Section~\ref{sec:human} uses the stratum-wise Clopper-Pearson construction behind the table's
intervals, taken one-sided; the difference uses each policy's construction at $\delta/2$. A prompt-clustered normal limit
for new prompts gives 15.7\% for the updated policy and 15.5\% for the pre-update model; the pre-update model's exact and
Wald-t limits are 20.4\% and 13.8\%.

\begin{table*}[t]
\centering
\scriptsize
\setlength{\tabcolsep}{2pt}
\resizebox{\ifdim\width>\textwidth\textwidth\else\width\fi}{!}{\begin{tabular}{@{}r r r r r p{3.2cm}@{}}
\toprule
\textbf{margin} & \textbf{chance of certifying} & \textbf{labels alone, unpaired} & \textbf{labels alone, paired by prompt} & \textbf{with the guard, pool rate} & \textbf{with the guard, new prompts} \\
\midrule
0.02 & 50\% & 1,226 & 416 & 167 & 339 \\
0.02 & 80\% & 2,801 & 950 & 380 & not within 490 prompts \\
0.03 & 80\% & 1,266 & 429 & 172 & 362 \\
0.05 & 80\% & 462 & 157 & 63 & 78 \\
\bottomrule
\end{tabular}
}
\caption{Prompt pairs needed to certify a margin on the difference of the two strict refusal rates at $\delta=0.05$, by
normal-type formula at the sheet's rates (difference $-0.0005$, pairing correlation 0.66, guard $\rho^2$ 0.6). The guard
columns assume a gain of 2.5; a design check measured about 1.4, so they are optimistic.}
\label{tab:cost}
\end{table*}

\paragraph{The prepared measurement.} One fresh response per policy on each of 490 prompts, 300 of the pairs forming the
labelling sheet, and 8 further responses per policy that only the guard reads, were drawn and scored; the analysis (an
exact betting bound, a bootstrap-t limit, and PPI++ with the guard for the pool rate and for new prompts) was fixed
before any label of this sample. On synthetic labels drawn from the guard's logit, the four limits' miss rates are at or
under 0.045 with the prompts fixed, and their mean certified margins 0.065, 0.037, 0.032 and 0.035: 300 pairs labelled
alone are expected to certify 2.4 points by the formula of Table~\ref{tab:cost} and, in this check (true difference $+0.006$), 3.7 points with the bootstrap-t limit and 6.5 with the exact bound, not the 2-point target.

\section{What Was Fixed in Advance}
\label{app:record}

A commit pushed to a remote before a result existed is a record of order we could not have written afterwards; where a
note and its results entered the repository together, the order rests on the note's own statement
(Table~\ref{tab:record}). Five of the 72 predictions say only that an exact bound is at or under its level; ten were
written after a dry run on earlier responses (Table~\ref{tab:d0}). Tables~\ref{tab:d1} to~\ref{tab:d4} list every
prediction and its outcome.

\begin{table*}[p]
\centering
\scriptsize
\begin{tabular}{@{}p{5.6cm} l p{9.2cm}@{}}
\toprule
\textbf{what was fixed} & \textbf{section} & \textbf{record} \\
\midrule
the first strata rule (ties kept together, cuts at quantiles, small strata merged) and the pre-flight's error criterion & App.~\ref{app:strata} & a design note committed on 2026-09-27, before the first results (the same day). Neither was kept: the rule was replaced and the criterion failed. The truth used for coverage, fixed in the same note, was changed as well (App.~\ref{app:strata}) \\
the expectations for the judge-assisted estimators (not the routing rule of Table~\ref{tab:routing}, which came out of the results) & App.~\ref{app:judge} & a note dated before the estimators were run, committed together with their results on 2026-09-30: no independent timestamp \\
the reading of a failure at step 200 & App.~\ref{app:judge} & a note dated while the scoring pass ran and before its output, committed together with the result on 2026-10-01: no independent timestamp \\
the hypotheses for the robot benchmark, against which the result on joint effort is marked exploratory & App.~\ref{app:bench} & a note written after the termination counts and 12 transcripts had been seen, committed together with the results on 2026-10-02: no independent timestamp \\
the two-way bounds: definitions, fallbacks, seed, sizes, reading and three expectations (Table~\ref{tab:d3}) & App.~\ref{app:bench} & a registration committed with its script on 2026-10-07, before the script was run on the benchmark's tables. The commit reached the remote only with its results, 42 minutes later, so the order rests on the local commit time \\
the confirmation pool: prompts, labels, class limits, checks, sizes, the GPU cap, eight predictions and their scoring code & App.~\ref{app:strata} & a registration committed and pushed with the scripts and the prompts on 2026-10-08, before any response was generated \\
the mid-rate pool: the mixture, labels, checks, sizes, the GPU cap, ten predictions and their scoring code & App.~\ref{app:strata} & a registration committed and pushed with the scripts, the prompts and the re-mix analysis on 2026-10-08, before any response on its prompts was generated. The predictions were written after the same analysis of a re-mix of the first two pools' responses, which loosened one (the bootstrap-t StratPPI limit may be over in one cell, by at most 0.01), added one (the Wilson-type limit within 0.01 of its level at a rate between 0.45 and 0.55) and set the floor of a third (the Wald-t limit's gain, at 80\% of the re-mix's smaller value); the registration names the first two as changes and gives the third's origin. Two more differences from the first registration were settled with the re-mix in view: the prediction on the cap is for the Wald-t limit's gain and not the Wilson-type limit's, and the Wilson-type limit with the sampled-pool term carries no prediction \\
the analysis of the human-label certificate & App.~\ref{app:human} & the rules and the analysis script, committed on 2026-10-04 before any label of this sample. They replaced a first version of the same day after the labels of App.~\ref{app:human} had shown a difference near zero. The treatment of a zero-variance resample in the approximate limits was amended once, on 2026-10-07, also before any label of this sample, with the design check unchanged. No label has been collected \\
the in-loop run and the costly-label test: environments or labels, designs, sizes, seeds, six predictions each and their scoring code & Apps.~\ref{app:pipeline}, \ref{app:strata} & a registration committed and pushed with the scripts on 2026-10-14, before any run; the costly-label script had printed data shapes and timings only \\
the exact bound and the suite-blocked resampling on AgentDojo: bounds, schemes, replications, seeds, six predictions and their scoring code & App.~\ref{app:bench} & committed and pushed with the script on 2026-10-14, before any run; a smoke test had printed timings and key names only. The run started 19 hours later than planned, unchanged, because the job meant to start it never did \\
two more seeds of the constrained run: seeds, steps, scoring wordings, item ids, the carried-bound computation and four predictions & App.~\ref{app:judge} & committed and pushed with the run script on 2026-10-14, before either seed was trained. One forced deviation: the first attempt at seed 1 crashed at step 12, and both seeds then ran without a registered memory-allocator setting, which changes memory allocation only; nothing of the crashed attempt is used. The registration reserved a retry of a seed cut short for an explicit decision; the seed was retried without one \\
the second model family: model, reward, pools, generation, checks, GPU caps, six predictions and their scoring code & App.~\ref{app:strata} & committed and pushed with the scripts and the prompts on 2026-10-14, before any training or any generation on its pool. The queue script's first pilot stopped on a fault in its own parser (an interpreter without the numerical library); the parser was fixed and the pilot rerun before the full run \\
exact stratified bounds and the chance of certifying: cells, arms, measures, eight predictions and their scoring code & App.~\ref{app:strata} & committed and pushed with the script on 2026-10-16, before any outcome on a registered cell was computed; the script had printed allocation sizes, two timings on made-up counts with their bound values, one unregistered toy cell, and Clopper-Pearson's exact miss on 3,000 made-up configurations of binomial sums. The step-spread sensitivity was written after the results and is not registered \\
AgentDojo against a fitted crossed model: the model, the fit, two schemes, the arms, five predictions and their scoring code & App.~\ref{app:bench} & committed and pushed with the script and the fitted parameters on 2026-10-16, before any bound was computed on a simulated table; the fit and a timing run had been seen \\
a second judge family and a second gold label: the model, items, four analyses, five predictions and their scoring code & App.~\ref{app:judge} & committed and pushed with the scripts on 2026-10-16, before any registered item was scored; a smoke test on four made-up items had been seen \\
the four limits of the human-label certificate from the existing sheet & App.~\ref{app:human} & a plan note written on 2026-10-19 before the limits were computed and after the sheet's rates were known; not committed before the computation: no independent timestamp, and not a test \\
the second model family trained on the guard's flag alone: the reward, the reused pool and strata, GPU caps, five predictions and their scoring code & App.~\ref{app:strata} & committed and pushed with the scripts on 2026-10-16, before any training step of the run; a 3-step pilot of the reward on training prompts had been seen (timing, memory, four training responses and the reward's spread). The registered reading took a failed check of the push's direction to mean the policy had not moved; the run moved it, and App.~\ref{app:strata} says so against that reading \\
\bottomrule
\end{tabular}

\caption{What was fixed before its result, and the record of the order.}
\label{tab:record}
\end{table*}

\begin{table*}[t]
\centering
\scriptsize
\begin{tabular}{@{}p{8.6cm} r r r@{}}
\toprule
 & \textbf{held} & \textbf{failed} & \textbf{untested} \\
\midrule
exact bounds at their level ($^\dagger$), registration pushed before the run & 5 & 0 & 0 \\
any other claim, registration pushed before the run & 37 & 17 & 0 \\
any other claim, written after a dry run & 9 & 0 & 1 \\
any other claim, local commit time only & 2 & 1 & 0 \\
all & 53 & 18 & 1 \\
\bottomrule
\end{tabular}

\caption{The registered predictions by content and record. $\dagger$: an exact bound at its level.}
\label{tab:d0}
\end{table*}

\begin{table*}[p]
\centering
\scriptsize
\begin{tabular}{@{}l p{7.4cm} p{7.4cm}@{}}
\toprule
 & \textbf{prediction, as registered} & \textbf{outcome} \\
\midrule
 & \emph{First registration: three pools, refusal rates 12\%, 82\% and 98\%} &  \\
P1 & large pool: the Wald-t limit and the bootstrap-t StratPPI limit are over in no mid-rate cell & refuted: the bootstrap-t limit is over in 1 of 8 cells (0.105 at $\delta$ 0.10, at 82\%); the Wald-t limit in 0 of 8 \\
P2a & large pool: the Wilson-type limit is over in no cell at a rate under 0.45 & kept: 0 of 4 cells at 12\% (0 of 12 with the two rare labels) \\
P2b & large pool: the Wilson-type limit is over in at least one of its four cells on every mid-rate label above 0.55 & kept: 4 of 4 at 82\% \\
P3 & the gain of the Wilson-type limit with a large pool is within 20\% of the pre-flight's prediction on every mid-rate label & kept: 1.53 and 1.62 against 1.82 at 12\%; 1.49 and 1.62 against 1.65 at 82\% \\
P4a & source: the Wilson-type limit with the sampled-pool term is over in no mid-rate cell at $\delta$ 0.05, and the Wald-t limit with the term in none at either $\delta$ & refuted: the Wilson-type limit is over in 1 of 4 cells (0.058 at 82\%, safety set 100); the Wald-t limit in 0 of 8 \\
P4b & source: the Wilson-type limit without the term and the bootstrap-t StratPPI limit are each over in at least a third of the mid-rate cells at $\delta$ 0.05 & kept: 3 of 4 and 4 of 4 \\
P4c & source: the gain of the Wilson-type limit with the term is within 15\% of the cap on every mid-rate label & kept: 1.32 and 1.22 against 1.35 and 1.24 at 12\%; 1.22 and 1.20 against 1.33 and 1.24 at 82\% \\
P5 & a rare label: the strata gain nothing (ESS of the Wilson-type limit between 0.9 and 1.15) and Clopper-Pearson on a random draw is over in no cell & kept: 1.01 to 1.11; 0 of 8 cells \\
 & \emph{Second registration, written after the re-mix: one mixed pool, refusal rate 46\%} &  \\
M1 & large pool: the Wald-t limit is over in no cell & kept: 0 of 4 (misses 0.026 and 0.036 at $\delta$ 0.05, 0.063 and 0.083 at 0.10) \\
M2 & large pool: the bootstrap-t StratPPI limit is over in at most one of four cells, and above its level by at most 0.01 in any & kept: 0 over, 3 unresolved, largest excess 0.001 \\
M3a & large pool: the Wilson-type limit is over in no cell at a rate under 0.45, and in at least one of four above 0.55 & not tested: the rate, 0.459, lies between \\
M3b & at a rate between 0.45 and 0.55, the miss of the Wilson-type limit is within 0.01 of its level in every cell & kept: 0.044 and 0.045 at $\delta$ 0.05, 0.096 at 0.10 \\
M4 & the gain of the Wilson-type limit with a large pool is within 20\% of the pre-flight's prediction & kept: 2.78 and 2.97 against 2.79 \\
M5 & the gain of the Wald-t limit with a large pool is at least 1.7 & kept: 2.09 and 2.56 \\
M6 & source: the Wald-t limit with the term is over in no cell & kept: 0 of 4 \\
M7 & source: the Wilson-type limit without the term and the bootstrap-t StratPPI limit are each over in at least one of the two cells at $\delta$ 0.05 & kept: 2 of 2 each (misses 0.079 to 0.120) \\
M8 & source: the gain of the Wald-t limit with the term is above 1 and within 15\% of the cap & kept: 1.49 and 1.34 against 1.64 and 1.44 \\
M9 & a rare label: ESS of the Wilson-type limit between 0.9 and 1.15, and Clopper-Pearson over in no cell & kept: 1.00 and 1.08; 0 of 4 cells \\
\bottomrule
\end{tabular}

\caption{Predictions on the strata's new pools. ``Large pool'': redrawn with replacement, 40,000 draws; ``source'': the
claim about the prompt source, 10,000 replications; a cell is one safety-set size at one $\delta$; gains are
$\mathrm{ESS}_{\text{truth}}$ at $\delta=0.05$. P1 to P5 were pushed before any response was generated; M1 to M9 were
pushed before generation on their prompts but written after a dry run on a re-mix of earlier responses.}
\label{tab:d1}
\end{table*}

\begin{table*}[p]
\centering
\scriptsize
\begin{tabular}{@{}l p{7.4cm} p{7.4cm}@{}}
\toprule
 & \textbf{prediction, as registered} & \textbf{outcome} \\
\midrule
 & \emph{The training loop (App.~\ref{app:pipeline}): synthetic bandit, 8 strata, four environments, 8,000 runs a cell} &  \\
Q1 & the Wald-t limit, for the pool's rate: over in at most 1 of its 8 cells & kept: 0 of 8 (2 unresolved) \\
Q2 & the Wald-t limit with the sampled-pool term, for the population's rate: over in at most 2 of its 8 cells & kept: 2 of 8 \\
Q3 & the same in the least heterogeneous environment at $\delta$ 0.10: over & kept: 0.107 \\
Q4 & the pooled Wilson bound on a random split, for the population's rate: over in at most 2 of its 8 cells & kept: 0 of 8 (2 unresolved) \\
Q5 & the $\mathrm{ESS}_{\mathrm{est}}$ of the Wald-t limit for the population's rate at $\delta$ 0.10 is at least 1.6 in the most heterogeneous environment & kept: 1.68 \\
Q6 & the same in the least heterogeneous environment is between 0.9 and 1.15 & kept: 1.10 \\
 & \emph{A costly label (App.~\ref{app:strata}): XSTest's human labels of five released models on 250 safe prompts} &  \\
Y1 & the Wald-t limit, labels at a rate of 10\% or more, safety set 100, with replacement: over in no cell & refuted: 3 of 20 cells, by at most 0.008 \\
Y2 & the Wald-t limit, mid-rate labels, without replacement: over in no cell & kept: 0 of 64 \\
Y3 & the Granite reference-rate stratifier, safety set 100, with replacement: the median gain of the Wald-t limit over the mid-rate labels is at least 1.1 & kept: 1.18 \\
Y4 & the same: no mid-rate label gains more than 2.0 & kept: at most 1.34 \\
Y5 & the stratifier from the other model families' completions gains more than the Granite one for at least 6 of the 8 mid-rate labels & refuted: 4 of 8 \\
Y6$^\dagger$ & Clopper-Pearson on a random draw, rare labels, with replacement: over in no cell & kept: no cell over \\
 & \emph{AgentDojo (App.~\ref{app:bench}): 28 pipelines, 4,000 resampled tables a pipeline and scheme} &  \\
Z1$^\dagger$ & the betting bound on user-task rates, against their average, is over for no pipeline with user tasks redrawn, across the table or within suites & kept: none under either \\
Z2$^\dagger$ & the pair-rate betting bound is over for no pipeline with user tasks redrawn within suites & kept: none \\
Z3 & the bound clustered by user task is over for no pipeline with user tasks redrawn within suites & kept: none \\
Z4 & the same bound is over for at least 15 of 28 pipelines with injection tasks redrawn within suites & refuted: 11 \\
Z5 & the quadrature bound is over for at most 2 of 28 pipelines under each scheme within suites & kept: none under any \\
Z6 & on the published tables the betting bound on user-task rates certifies 5\% for at most one pipeline, and the pair-rate one for none & kept: one (0.048) and none \\
 & \emph{A second model family (App.~\ref{app:strata}): Qwen2.5-1.5B-Instruct, a new pool of 400 prompts} &  \\
V1 & at step 200 $\rho$ on the refusal label is under 0.8 & refuted: 0.84 \\
V2 & at step 200 the refusal rate rose on the categories the reward pushed up and fell on those it pushed down & refuted: 16.9\% to 14.8\%, and 22.5\% to 23.5\% \\
V3 & at step 200, at a refusal rate between 5\% and 95\%, the Wald-t limit with a large pool is over in no cell & kept: 0 of 4 at 19\% \\
V4 & the Wald-t limit's gain with a large pool (safety set 200) is lower at step 200 than at step 100 & kept: 1.77 against 1.84 \\
V5 & at step 200 the gain of the Wilson-type limit with a large pool (safety set 200) is within 20\% of the pre-flight at the measured moderators & kept: 1.70 against 1.73 \\
V6 & at step 200 the Wald-t limit's gain with a large pool (safety set 200) is between 0.9 and 1.6 & refuted: 1.77 \\
 & \emph{Two more seeds of the constrained run (App.~\ref{app:judge}), step 200 against step 0} &  \\
S1 & in at least one of the two new seeds the carried bound is over its level for at least one of the six wordings & kept \\
S2 & the same in both new seeds & kept: 3 and 2 of 6 \\
S3 & in both new seeds the canonical wording's recall on the guard's refusals is lower at step 200 than at step 0 & kept: 0.21 to 0.07 and to 0.05 \\
S4 & over the three seeds together, at least 6 of the 18 carried bounds are over & kept: 8 of 18 \\
\bottomrule
\end{tabular}

\caption{Predictions registered on one day, each pushed before any of its runs: the training loop (Q), a costly human label (Y),
AgentDojo blocked by suite (Z), a second model family (V), two further seeds of the targeted update (S).}
\label{tab:d2}
\end{table*}

\begin{table*}[t]
\centering
\scriptsize
\begin{tabular}{@{}l p{7.4cm} p{7.4cm}@{}}
\toprule
 & \textbf{expectation, as registered} & \textbf{outcome} \\
\midrule
T1 & the multiway variance with a t quantile is over for some pipelines & kept: 26 of 28 \\
T2 & the pigeonhole bootstrap-t is over for fewer pipelines than it & kept: none \\
T3 & the quadrature bound is over for none & refuted: one (0.058) \\
\bottomrule
\end{tabular}

\caption{Expectations for the two-way bounds on AgentDojo, scheme (c); the order rests on a local commit time.}
\label{tab:d3}
\end{table*}

\begin{table*}[p]
\centering
\scriptsize
\begin{tabular}{@{}l p{7.4cm} p{7.4cm}@{}}
\toprule
 & \textbf{prediction, as registered} & \textbf{outcome} \\
\midrule
 & \emph{Exact stratified bounds and the chance of certifying (App.~\ref{app:strata}): 116 cells, $\delta$ 0.05 and 0.10 for X1 and X2, $\delta$ 0.05 for X3 to X8} &  \\
X1$^\dagger$ & neither exact stratified bound is over in any cell & kept: 0 of 464 \\
X2 & the betting bound's width gain against Clopper-Pearson is below 1 in every cell & kept: at most 0.53 \\
X3 & the betting bound's power ESS at a pass probability of 0.8 is below 1 in every real and costly cell at a rate of 5\% or more & kept: at most 0.28 \\
X4 & the stratified Clopper-Pearson bound's power ESS at 0.8, real and costly cells in range: lowest at least 0.95, median at least 1.10 & refuted: 0.27 and 0.92 \\
X5 & the same at a pass probability of 0.5 lies between 0.85 and 1.15 in every one & refuted: 0.07 to 1.09 \\
X6 & the same at 0.8 on rare real and costly cells lies between 0.8 and 1.2 & refuted: 0.33 to 1.01 \\
X7 & real cells in range: the Wald-t limit's width gain against Clopper-Pearson is at least its gain against the pooled Wilson bound in every one & kept: 36 of 36 \\
X8 & the same cells: the Wald-t limit's power ESS at 0.8 is above the stratified Clopper-Pearson bound's in every one & refuted: 35 of 36 \\
 & \emph{AgentDojo against a fitted crossed model (App.~\ref{app:bench}): 28 pipelines, 4,000 tables a pipeline and scheme} &  \\
W1$^\dagger$ & injection tasks fixed: the betting bound on user-task rates (against their average) and the pair-rate betting bound are over for no pipeline & kept: none and none \\
W2 & the same scheme: the bound clustered by user task is over for at most 2 & kept: none \\
W3 & both kinds of task drawn: the bound clustered by user task is over for at least 8 & kept: 18 \\
W4 & the same scheme: the quadrature bound is over for at most 2 & kept: none \\
W5 & the same scheme: the betting bound on user-task rates is over for at least 2 & kept: 5 \\
 & \emph{A second judge family and a second gold label (App.~\ref{app:judge}): WildGuard} &  \\
J1 & WildGuard as the judge, the guard as gold: the carried bound at step 200 is over for at least 1 of the 3 seeds & refuted: none (misses 0.000) \\
J2 & the same: WildGuard's recall on the guard's refusals is lower at step 200 than at step 0 in all 3 seeds & refuted: higher in all 3 (0.88 to 0.94, 0.94 and 0.91) \\
J3 & WildGuard as gold, the compiled judge's canonical wording as the judge: the carried bound at step 200 is over for at most 1 of the 3 seeds & refuted: all 3 (0.84, 0.46, 0.66) \\
J4 & on the refusal sheet at step 200, WildGuard's design-weighted accuracy against the human strict label is higher than the guard's & kept: 0.930 against 0.929 \\
J5 & among the guard's flags, WildGuard flags a smaller share of responses with the surface pattern than without (Fisher p < 0.05) & refuted: 95\% against 92\% (p 0.064) \\
 & \emph{The second model family trained on the guard's flag alone (App.~\ref{app:strata}): the same pool and strata} &  \\
U1 & at step 200 refusal on the categories pushed up is at least 15 points above the reference, and on those pushed down at least 10 points below it & refuted: 16.9\% to 67.4\%, and 22.5\% to 48.7\% \\
U2 & at step 200 $\rho$ on the refusal label is under 0.7 & kept: 0.46 \\
U3 & at step 200, at a refusal rate between 5\% and 95\%, the Wald-t limit with a large pool is over in no cell & kept: 0 of 4 at 58\% \\
U4 & at step 200 the Wald-t limit's gain with a large pool (safety set 200, $\delta$ 0.05) is below 1.3 & kept: 1.08 \\
U5 & at step 200 the gain of the Wilson-type limit with a large pool (safety set 200) is within 25\% of the pre-flight at the measured moderators & kept: 1.22 against 1.14 \\
\bottomrule
\end{tabular}

\caption{Predictions registered together, before any of their runs: exact stratified bounds (X), AgentDojo against a
fitted crossed model (W), a second judge family (J), and the update that reordered the
prompts (U). A power ESS is the number of random labels Clopper-Pearson needs to certify as often, over the labels the
bound uses.}
\label{tab:d4}
\end{table*}

\section{Reproducibility}

Every number comes from scripts and result files that will be released with the paper. The resampling studies ran on
the CPU of one workstation (AMD Ryzen 7 7800X3D, 8 cores, 24\,GB of memory available to Ubuntu under WSL2); the
policies were trained and sampled, and the guard and judges run (4-bit where stated), on one NVIDIA GeForce RTX 4070
SUPER with 12\,GB, with Python 3.12, PyTorch 2.13, Transformers 5.16, TRL 1.12, PEFT 0.20 and bitsandbytes 0.50. Each
resampling study fixes its random seeds in its script. Human
labels were collected under a written guideline that follows XSTest's classes. Draw counts per cell are given with each
table.
